\section{Tautologi, Kontradiksi, dan Kontingensi}

Klasifikasi proposisi majemuk berdasarkan nilai kebenarannya merupakan fondasi penting sebelum membahas ekuivalensi logis dan inferensi formal \cite{rosen2019}. Secara umum, proposisi majemuk dikelompokkan menjadi tiga jenis, yaitu tautologi, kontradiksi, dan kontingensi.

\subsection{Tautologi}

\begin{definisi}[Tautologi]
\logic{Tautologi} adalah proposisi majemuk yang \textbf{selalu bernilai benar (True)} untuk setiap kombinasi nilai kebenaran proposisi atomik penyusunnya.
\end{definisi}

Dalam konteks pemrograman atau basis data, ekspresi yang selalu benar tidak memberikan penyaringan kondisi yang efektif karena hasil evaluasinya selalu meloloskan.

\begin{contoh}
Contoh tautologi paling sederhana adalah \textbf{hukum tiada jalan tengah}: $p \vee \neg p$. \\
Pernyataan ini selalu bernilai benar, apa pun nilai kebenaran $p$.
\end{contoh}

\subsection{Kontradiksi}

\begin{definisi}[Kontradiksi]
\logic{Kontradiksi} adalah proposisi majemuk yang \textbf{selalu bernilai salah (False)} untuk setiap kombinasi nilai kebenaran proposisi atomik penyusunnya.
\end{definisi}

Dalam penerapan komputasi, kondisi kontradiktif dapat menyebabkan cabang program tidak pernah dieksekusi (\textit{dead code}).

\begin{contoh}
Contoh kontradiksi paling sederhana adalah \textbf{hukum non-kontradiksi}: $p \wedge \neg p$. \\
Pernyataan ini selalu bernilai salah karena $p$ dan $\neg p$ tidak mungkin benar pada saat yang sama.
\end{contoh}

\subsection{Kontingensi}

\begin{definisi}[Kontingensi]
\logic{Kontingensi} adalah proposisi majemuk yang nilai kebenarannya \textbf{bergantung} pada kombinasi nilai proposisi atomik penyusunnya; pada sebagian kombinasi bernilai benar, dan pada kombinasi lain bernilai salah.
\end{definisi}

Sebagian besar ekspresi logika dalam praktik, seperti $p \rightarrow q$ atau $(p \vee q) \wedge \neg r$, termasuk kontingensi. Oleh karena itu, evaluasi umumnya memerlukan tabel kebenaran atau manipulasi aljabar yang tepat.
